% !TEX root = userguide.tex
\def\headdate{30th Sep 2021}

\section{The diffuse-interface model}
\label{sec:examples_dim}

In this section an example using the diffuse-interface model will be given
after the theory in the next section.

\subsection{Weak form of the diffuse-interface model. Galerkin \fem}

We consider the diffuse-interface model in a region $\Omega$, with the flow
equations according to Stokes flow:
\begin{alignat}{3}
-&\nabla\cdot(2\eta \ten D)+\nabla g_0&&=\vek f+\nabla \cdot \ten    
\tau_\text{c}, &&\qquad\text{in }\Omega \label{eq:di1}\\
&\nabla\cdot\vek u &&=0 &&\qquad\text{in }\Omega \label{eq:di2}
\end{alignat}
where $\ten D=(\nabla\vek u+\nabla\vek u^T)/2$, $g_0$ is a pressure term, and 
$\ten \tau_\text{c}$ is the 
capillary stress tensor given by
\[
\ten \tau_\text{c} = \rho \kappa (|\nabla c|^2 \ten I - \nabla c \nabla c)
\]
where $c$ is the concentration, $\kappa$ is a parameter related to the thickness
of the interface,  and $\rho$ is the density of the fluid (the inertia
terms in the momentum balance, which also contain $\rho$, are assumed to be 
negligible). Furthermore, an isotropic stress 
$\rho \kappa |\nabla c|^2 \ten I$ was added to ensure 
that the capillary stress is always parallel to the interface. We will call 
this form of the momentum balance, where the capillary stress is used directly, 
the \emph{stress form}. The divergence 
of the capillary stress tensor $\nabla \cdot \ten \tau_\text{c}$ can be 
rewritten to a body force $\rho \mu \nabla c$, or alternatively $-\rho c\nabla 
\mu$, plus an isotropic term that can be absorbed in the pressure. We will call 
these two forms of the 
flow equations \emph{potential form 1} (with pressure $g_1$) and 
\emph{potential form 
	2} (with pressure $g_2$), 
respectively \cite{Jaensson2015}. Note, that 
potential form 1 is employed in most of the 
diffuse-interface examples in \tfem.

The equations for the
concentration $c$ and chemical potential $\mu$ are given by:
\begin{alignat}{3}
&\pderiv{c}{t}+\vek u\cdot\nabla c&-M\nabla^2\mu=0&&\qquad\text{in }\Omega
\label{eq:di3}\\
\alpha c&-\beta c^3+\kappa\nabla^2 c&+\mu\,\,\,\,=0&&\qquad\text{in }\Omega
\label{eq:di4}
\end{alignat}
where $M$ is the mobility, and $\alpha$ and $\beta$ are diffuse-interface
parameters related to the shape of the double-well function.
For a planar interface (varying only in $z$-direction) which is in 
equilibrium ($\mu = 0$), the concentration
profile can be determined:
\begin{equation}
c(z) = c_\text{B} \tanh\left(\frac{z}{\sqrt{2} \xi}\right)
\end{equation}
where the bulk values are given by $c_\text{B}=\pm \sqrt{\alpha/\beta}$ and 
$\xi = \sqrt{\kappa / \alpha} $. In this case, the interfacial tension $\gamma$ 
is given by 
\begin{equation}
\gamma = \frac{2 \sqrt{2}}{3} \frac{\rho \kappa c_\text{B}^2}{\xi}
\label{eq:di_surface_tension}
\end{equation}

The boundary conditions for the
flow problem are the same as in Sec.~\ref{sec:stokesweak}.
The boundary conditions for $c$ and $\mu$ are assumed to be
$\plderiv{c}{n}=0$ and $\plderiv{\mu}{n}=0$ on $\Gamma$, which correspond to a 
90 degree contact angle and no influx of $c$, respectively. For a contact angle 
different from 90 degrees see 
Sec.~\ref{sec:contactangle}.                                                    


The weak form of the flow equations using the stress form is similar to 
Sec.~\ref{sec:stokesweak}:
find $\vek u$ and $g_0$ such that
\begin{alignat*}{3}
&\big(\ten D_v,2\eta\ten D\big)-(\nabla\cdot\vek v,g_0)&&=(\vek v,
\vek t_{\text{N}})_{\Gamma_{\text{N}}}+(\vek v,\vek f)-(\ten D_v,\ten 
\tau_\text{c})&\qquad&\text{for all }\vek v\\
&\,(q,\nabla\cdot\vek u)&&=0&\qquad&\text{for all }q
\end{alignat*}
using appropriate spaces for $\vek u$, $g_0$, $\vek v$ and $q$. Using potential 
form 1, we obtain the weak form: 
find $\vek u$ and $g_1$ such that
\begin{alignat*}{3}
&\big(\ten D_v,2\eta\ten D\big)-(\nabla\cdot\vek v,g_1)&&=
(\vek v,\vek f)+(\vek v,\rho \mu \nabla c)&\qquad&\text{for all }\vek v\\
&\,(q,\nabla\cdot\vek u)&& =0&\qquad&\text{for all }q
\end{alignat*}
And similarly for potential form 2:
find $\vek u$ and $g_2$ such that
\begin{alignat*}{3}
&\big(\ten D_v,2\eta\ten D\big)-(\nabla\cdot\vek v,g_2)&&= (\vek v,\vek f)   
-(\vek v,\rho c \nabla \mu)&\qquad&\text{for all }\vek v\\
&\,(q,\nabla\cdot\vek u)&& =0 &\qquad&\text{for all }q
\end{alignat*}

Note that the pressure fields $g_0$, $g_1$ and $g_2$ are defined differently 
and therefore cannot be compared directly 
\cite{Jaensson2015}. 

In these equations
$\ten D_v$ is defined by
\[
\ten D_v=(\nabla\vek v+\nabla\vek v^T)/2.
\]
It is important to note that adding Neumann (traction) boundary conditions for 
the momentum balance 
using one of the potential forms is not straightforward, as was shown in 
\cite{Jaensson2015}. Therefore, Neumann (traction) boundary conditions are left 
out of the weak forms of the momentum balance using the potential forms.

The weak form of the $c$ and $\mu$ equations can easily be derived and reads:
find $c$ and $\mu$ such that
\begin{alignat*}{3}
\big(r,\pderiv{c}{t}+\vek u\cdot\nabla c\big)+M(\nabla r,\nabla \mu)
&=0 &&\qquad&\text{for all } r\\
(s, \alpha c-\beta c^3)-\kappa(\nabla s,\nabla c) +(s,\mu) &=0
&&\qquad&\text{for all } s
\end{alignat*}
using appropriate spaces for $c$, $\mu$, $r$ and $s$.


The weak form can be used to obtain an approximate solution using the Galerkin 
\fem\
technique. For this we divide the region into elements:
\[
\Omega=\cup_e\Omega_e, \quad\Omega_{e}\cap\Omega_{f}=\emptyset\quad
\text{for }e\neq f
\]
and interpolate $\vek u$, $g_0$, $c$ and $\mu$
(and similarly $\vek v$, $g_1$, $g_2$, $q$, $r$ and $s$) by polynomials on each 
element:
\[
\begin{split}
\vek u_h&=\tsum_k \phi_k(\vek x)\vek u_k=\col\phi^T(\vek x)\col{\vek u}\\
[g_0]_h&=\tsum_k \psi_k(\vek x) [g_0]_k=\col\psi^T(\vek x)\col{g_0}\\
c_h&=\tsum_k \phi_k(\vek x) c_k=\col\phi^T(\vek x)\col{ c}\\
\mu_h&=\tsum_k \phi_k(\vek x) \mu_k=\col\phi^T(\vek x)\col{ \mu}
\end{split}
\]
where $\col\phi(\vek x)$ and $\col\psi(\vek x)$
are global shape functions. Note, that $c$ and $\mu$ are interpolated using the
same shape function as the velocity vector.

\subsection{Explicit/implicit time discretization 
	\label{sec:di_expl}}
At the initial time $t_0$ we solve the flow equations to obtain an initial flow 
field $\vek u^0=\vek u(t_0)$. This can be done with or without
the right-hand side term $\nabla \cdot \ten \tau_\text{c}$ (which can be 
replaced by $\rho\mu\nabla c$ or $-\rho c\nabla \mu$ when one of the potential 
forms is used). Due to the time derivative of $c$, the initial field 
$c^0=c(t_0)$ should be specified.
Then we do a time
stepping where we first solve the equations for 
$c^{n+1}$ and $\mu^{n+1}$ using the values of $\vek u^n$ and $c^n$ at the
current time $t^n$. Using an implicit Euler scheme, we obtain
\begin{alignat*}{3}
&\big(r,\frac{c^{n+1}-c^n}{\Delta t}+\vek u^n\cdot\nabla c^{n+1}\big)&
+M(\nabla r,\nabla \mu^{n+1})=0 &&\qquad&\text{for all } r\\
(s,\alpha c^{n+1}&-\beta [c^{n+1}]^3)-\kappa(\nabla s,\nabla c^{n+1})   
&+(s,\mu^{n+1})=0 &&\qquad&\text{for all } s 
\end{alignat*}
where $\Delta t$ is the time step. This is a non-linear equation, which we
solve by Picard iteration:
\begin{alignat}{3}
&\big(r,\frac{c_{i+1}-c^n}{\Delta t}+\vek u^n\cdot\nabla   
c_{i+1}\big)&+M(\nabla r,\nabla \mu_{i+1})=0 &&\qquad&\text{for all } r\\
\big(s,(\alpha &-\beta c_{i}^2)c_{i+1}\big)-\kappa(\nabla s,\nabla c_{i+1}) 
&+(s,\mu_{i+1})=0 &&\qquad&\text{for all } s \label{eq:di_picard}
\end{alignat}
or by Newton-Raphson iteration:
\begin{alignat}{3}
&\big(r,\frac{c_{i+1}-c^n}{\Delta t}+\vek u^n\cdot\nabla c_{i+1}\big)&+ 
M(\nabla r,\nabla \mu_{i+1})=0 &&\qquad&\text{for all } r\\
\big(s,(\alpha &-3\beta c_{i}^2)c_{i+1}\big)-\kappa(\nabla s,\nabla c_{i+1})
&+(s,\mu_{i+1})= -\big(s,2\beta c_{i}^3 \big) &&\qquad&\text{for all } s
\label{eq:di_newton}
\end{alignat}
for $i=0,1,\ldots$ with $c_0=c^n$, until convergence.
Then we solve the flow equations to obtain $\vek u^{n+1}$ and $p^{n+1}$ using
$c^{n+1}$ and $\mu^{n+1}$.

Replacing the implicit Euler scheme with a Gear (BDF2) scheme,
we obtain second-order time discretization:
\begin{alignat*}{3}
&\big(r,\frac{\frac32c^{n+1}-2c^n+\frac12c^{n-1}}{\Delta t}+
\hat{\vek u}^{n+1}\cdot\nabla c^{n+1}\big)&+M(\nabla r,\nabla \mu^{n+1})
=0 &&\qquad&\text{for all } r\\
&(s, \alpha c^{n+1}-\beta [c^{n+1}]^3)-\kappa(\nabla s,\nabla c^{n+1})
&+(s,\mu^{n+1})=0 &&\qquad&\text{for all } s
\end{alignat*}
where $\hat{\vek u}^{n+1}=2\vek u^n-\vek u^{n-1}$ is a prediction of the
velocity vector field $\vek u$ at $t^{n+1}$.

Remarks:
\begin{enumerate}
	\item The first-order discretization has been the base of the original
	element routine\\ \texttt{diffuse\_interface\_elem}.
	\item The second-order discretization has been the implemented in the
	element routine \texttt{diffuse\_interface\_elem1}, along with the
	extensions for ALE and tALE (see Sec.~\ref{sec:ale}).
	\item In a scheme where the flow is computed first we can replace 
	$\hat{\vek u}^{n+1}$ with the computed $\vek u^{n+1}$, however 
	we need a predicted $\hat{c}^{n+1}=2c^n-c^{n-1}$ in order to obtain
	second-order accuracy in time.
	\item Picard iteration is used in  \texttt{diffuse\_interface\_elem}.
	If \texttt{diffuse\_interface\_elem1} is used, both Picard and
	Newton-Raphson iteration is available. 
\end{enumerate}

\subsection{Planar Couette flow}
\label{sec:couettedi}

We consider a planar Couette problem on a unit square.
On the upper boundary a tangential
velocity of 0.5 and a normal velocity of zero is imposed.
On the lower boundary a tangential
velocity of -0.5 and a normal velocity of zero is imposed.
In horizontal direction periodical boundary conditions are assumed.
In the lower left corner the pressure value is set to zero (Dirichlet).

At the start of the flow the concentration $c$ is set to zero
plus a small random perturbation.

The example file is \texttt{diffuse\_interface1.f90}
and can be obtained by typing at the
command line:
\medskip
\begin{quote}
   \texttt{getexample diffuse\_interface}
\end{quote}
The post-processing part is done by the separate programs
\texttt{diffuse\_interface\_post.f90} and \texttt{streamfunction.f90}.

The example is rather straightforward and can be studied by examining the
source files. The periodical boundary conditions are imposed using
constraints (see Chap.~\ref{ch:constraints}).
The add-on `diffuse interface' (see Sec.~\ref{sec:diffuse_interface})
is used for the building of the system. 

The example can be run by
\begin{quote}
   \texttt{diffuse\_interface1 < input\_diffuse\_interface1.txt}
\end{quote}
which might take a while.
The results are written to a file and can be turned into viewable fig files by
\begin{quote}
   \texttt{diffuse\_interface\_post}\\
   \texttt{streamfunction}
\end{quote}

After running the example
the results in Fig.~\ref{fig:couette1di} are obtained.
\begin{figure}[htbp!]
   \begin{center}
   \includegraphics[scale=0.5]{c_color_fill}
   \end{center}
   \caption{Concentration $c$ at time 1 
            in a planar Couette flow with parameters $M=0.02$,
            $\kappa=4\times10^{-4}$ and $\rho=50$. All others parameters
            ($\eta$, $\alpha$, $\beta$) are set to 1.
            Mesh: 80$\times$80 elements. Time step $\Delta t=1\times10^{-2}$.}
    \label{fig:couette1di}
\end{figure}

A similar example is \texttt{diffuse\_interface11.f90}, but now using spectral
elements. 
The example can be run by
\begin{quote}
   \texttt{diffuse\_interface11 < input\_diffuse\_interface11.txt}
\end{quote}
which might take a while.
The results are written to a file and can be turned into viewable fig files by
\begin{quote}
   \texttt{diffuse\_interface\_post11}\\
   \texttt{streamfunction11}
\end{quote}

Another similar example is \texttt{diffuse\_interface12.f90}, but now
using second-order time discretization.
The example can be run by
\begin{quote}
   \texttt{diffuse\_interface12 < input\_diffuse\_interface12.txt}
\end{quote}
which might take a while.
The results are written to a file and can be turned into viewable fig files by
\begin{quote}
   \texttt{diffuse\_interface\_post}\\
   \texttt{streamfunction}
\end{quote}

\subsection{Biperiodic shear flow. Lees-Edwards boundary conditions}
\label{sec:couettedi5}

We consider a biperiodic shear flow problem on a unit square.
The upper and lower boundary are coupled using Lees-Edwards boundary
conditions (LEbc) with velocity difference of 1 between the upper and lower
boundary of the domain.
In horizontal direction periodical boundary conditions are assumed.
Halfway on the left boundary the velocity is set to zero (Dirichlet).
In the lower left corner the pressure value is set to zero (Dirichlet).

At the start of the flow the concentration $c$ is set to zero
plus a small random perturbation.

The example file is \texttt{diffuse\_interface5.f90}
and can be obtained by typing at the
command line:
\medskip
\begin{quote}
   \texttt{getexample diffuse\_interface}
\end{quote}
The post-processing part is done by the separate programs
\texttt{diffuse\_interface\_post.f90} and \texttt{streamfunction.f90}.

The example is rather straightforward and can be studied by examining the
source files. The (periodical) boundary conditions are imposed using
weak constraints on objects (see Chap.~\ref{ch:constraints}).
Some notes on the source:
\begin{listing}{481}
!   shift mapped coordinates forward in time for Lees-Edwards conditions

    time = time + deltat

    call find_refcoor_objects ( mesh, object1=4, mapcoor=mapcoor )
\end{listing}
Here the intersections of the coordinates of object 4 are recomputed every time
step. The function \texttt{mapcoor} maps the coordinates according to
Lees-Edwards boundary conditions before the intersection is computed.

NOTE: the weak constraint integrals need to be computed
exact to get the best results (see \cite{Hwang04}).
In this example we achieve this in the following way. 
The number of elements along the object is 80. The integration rule is
two-point Gauss on 5 subintervals. This means that if we take the time step to
be a multiple of $1/(5*80)=2.5\times10^{-3}$ the integrals are computed 
exactly.
A disadvantage of this approach is that if the shear rate is not constant the
integrals cannot be integrated exactly. In that case much more integration
points need to used or, even better, the integration should be performed
analytically (like in \cite{Hwang04}) or numerically with a split interval
(like in \cite{Anderson06}).

The example can be run by
\begin{quote}
   \texttt{diffuse\_interface5 < input\_diffuse\_interface5.txt}
\end{quote}
which might take a while.
The results are written to a file and can be turned into viewable fig files by
\begin{quote}
   \texttt{diffuse\_interface\_post}
\end{quote}

After running the example
the results in Fig.~\ref{fig:couette5di} are obtained.
\begin{figure}[htbp!]
   \begin{center}
   \includegraphics{c_color_fill5}
   \end{center}
   \caption{Concentration $c$ at time 1 
            in a biperiodic (with LEbc)
            flow with parameters $M=0.02$,
            $\kappa=4\times10^{-4}$ and $\rho=50$. All others parameters
            ($\eta$, $\alpha$, $\beta$) are set to 1.
            Mesh: 80$\times$80 elements. Time step $\Delta t=1\times10^{-2}$.}
    \label{fig:couette5di}
\end{figure}
\subsection{More diffuse-interface examples}
\label{sec:diexamples}

(These examples have been developed by Patrick Anderson of the TU/e).

One of the nice features of the diffuse-interface model is that it can easily 
be used to study a variety of different problems by changing
the concentration $c$. The example file is \texttt{diffuse\_interface3.f90}
and is obtained by typing at the
command line:
\medskip
\begin{quote}
   \texttt{getexample diffuse\_interface}
\end{quote}
The post-processing part is again done by the separate programs
\texttt{diffuse\_interface\_post.f90} and \texttt{streamfunction.f90}.

Below, a part of the source file is given defining the initial concentration 
vector such that retraction of a non-circular drop is studied:
\begin{listing}{250}
! create concentration vector
!
!   funcnr = 1 : unstable stratified layers
!   funcnr = 2 : retraction of non-circular drop
!   funcnr = 3 : spinodal decomposition
!   funcnr = 4 : a supercritical and a subcritical nucleus

  call fill_sysvector ( mesh, problemdi, soldin, degfd=1, &
       node1=1, node2=mesh%nnodes, func=difunc, funcnr=2)
\end{listing}

Again this example is rather straightforward and can be studied by examining 
the source files. In this example bi-periodic boundary conditions are also 
imposed  using constraints (see Chap.~\ref{ch:constraints}), but now coupling 
the top and bottom curve and the two side curves.
The add-on `diffuse interface' (see Sec.~\ref{sec:diffuse_interface})
is used for the building of the system.

The example can be run by
\begin{quote}
   \texttt{diffuse\_interface3 < input\_diffuse\_interface3.txt}
\end{quote}
which might take a while depending on the exact conditions.
The results are written to a file and can be turned into viewable fig files by
\begin{quote}
   \texttt{diffuse\_interface\_post}
\end{quote}

\begin{figure}[htbp!]
   \begin{center}
   \resizebox{0.2\textwidth}{!}{
     \includegraphics*{dropretract0}
   }
   \resizebox{0.2\textwidth}{!}{
     \includegraphics*{dropretract10}
   }
   \resizebox{0.2\textwidth}{!}{
     \includegraphics*{dropretract20}
   }
   \resizebox{0.2\textwidth}{!}{
     \includegraphics*{dropretract50}
   }
   \end{center}
   \caption{Concentration $c$ at different times 
            in a bi-periodic flow domain with parameters $M=0.04$,
            $\kappa=16\times10^{-4}$ and $\rho=250$. All others parameters
            ($\eta$, $\alpha$, $\beta$) are set to 1.
            Mesh: 40$\times$40 elements. Time step $\Delta t=1\times10^{-2}$.}
    \label{fig:diexample3}
\end{figure}


\begin{figure}[htbp!]
   \begin{center}
   \resizebox{0.2\textwidth}{!}{
     \includegraphics*{bingrowth0}
   }
   \resizebox{0.2\textwidth}{!}{
     \includegraphics*{bingrowth05}
   }
   \resizebox{0.2\textwidth}{!}{
     \includegraphics*{bingrowth050}
   }
   \resizebox{0.2\textwidth}{!}{
     \includegraphics*{bingrowth100}
   }
   \end{center}
   \caption{A supercritical and a subcritical nucleus; 
            concentration $c$ at different times
            in a bi-periodic flow domain with parameters $M=0.04$,
            $\kappa=16\times10^{-4}$ and $\rho=250$. All others parameters
            ($\eta$, $\alpha$, $\beta$) are set to 1.
            Mesh: 40$\times$40 elements. Time step $\Delta t=1\times10^{-2}$.}
    \label{fig:diexample3a}
\end{figure}


A three-dimensional example of the diffuse-interface model in a brick
with no external flow, homogeneous essential boundary conditions for 
velocities 
can be run by
\begin{quote}
   \texttt{diffuse\_interface4 < input\_diffuse\_interface4.txt}
\end{quote}
which really take a while depending on the exact conditions.
This example generates VTK data files (see
Sec.~\ref{sec:vtk}) for plotting the results.

An axisymmetric example of a drop deforming in extensional flow is provided 
in a sixth example, i.e.
\texttt{diffuse\_interface6.f90} where the constant
\begin{quote}
    coorsys = 1       ! axisymmetric coordinate system
\end{quote}
has been defined. After 250 time steps the boundary conditions are set to 
zero, 
the drop relaxes and breaks up into multiple drops. This example uses the 
commercial package Tecplot (see Sec.~\ref{sec:tec}) for post processing at 
each 
tenth time step. A one-dimensional example is provided in 
\texttt{diffuse\_interface7.f90} where only 
the equations for $c$ and $\mu$ are solved to study spinodal decomposition in 
great 
detail. Example \texttt{diffuse\_interface8.f90} can be used to study the 
effect 
of hydrodynamic forces on growth rates in chaotic flow. The program 
demonstrates how 
to apply a time-periodic source term for the momentum equation.

\subsection{The ternary diffuse-interface model}

(This module have been developed by Patrick Anderson of the TU/e).

Also a ternary version of the diffuse-interface model is available in \tfem.
The main difference compared to the two-component Cahn-Hilliard model is that 
now
a third phase is present; thus the free energy is now a function of $c_1$ and 
$c_2$.
Given the ternary free energy
\begin{equation}
f_0 = \frac{1}{4} 
\left(
c_1^2 c_2^2 + (c_1^2+c_2^2)\times(1-c_1-c_2)^2 -c_1c_2 (1-c_1-c_2)
\right)
\end{equation}
the derivatives read
\begin{align*}
\frac{\partial f_0}{\partial c_1} &=
0.5c_1 
-1.5c_1^2+c_1^3-0.25c_2-0.5c_1c_2+1.5c_1^2c_2-0.25c_2^2+1.5c_1c_2^2+0.5c_2^3\\
&= ( 0.5 -1.5c_1+c_1^2 - 0.5c_2+1.5c_1c_2
+1.5c_2^2)c_1 + (-0.25 -0.25c_2 +0.5c_2^2)c_2\\
\frac{\partial f_0}{\partial c_2} &=
( 0.5 -1.5c_2+c_2^2 - 0.5c_1+1.5c_1c_2
+1.5c_1^2)c_2 + (-0.25 -0.25c_1 +0.5c_1^2)c_1
\end{align*}
The system of equations for $\mu_1,\mu_2, c_1, c_2,\vec u, p$, derived by  
Mauri, Molin and Anderson, now reads
\begin{align}
\mu_1 &= \frac{\partial f_0}{\partial c_1}-{\epsilon^2}\nabla^2 c_1 - 
\frac{\epsilon^2}{2} \nabla^2 c_2 \\
\mu_2 &= \frac{\partial f_0}{\partial c_2} -{\epsilon^2}\nabla^2 c_2 - 
\frac{\epsilon^2}{2} \nabla^2 c_1 \\
\frac{\partial c_1}{\partial t} +{\vec u} \cdot \nabla c_1 &=  \nabla  \cdot M 
c_1 c_2 \nabla \mu_2 +
 \nabla  \cdot M c_1 (1-c_2) \nabla i\mu_1\\
\frac{\partial c_2}{\partial t} +{\vec u} \cdot \nabla c_2 &=  \nabla  \cdot 
M  
c_1 c_2 \nabla \mu_1 +
 \nabla  \cdot M c_1 (1-c_2) \nabla \mu_2
\end{align}
solved in combination with the mass and momentum balance
\begin{align}
0&= -\nabla p + \eta\nabla^2 {\vec u} + \rho (\mu_1\nabla c_1 + \mu_2\nabla 
c_2)\\
\nabla \cdot {\vec u}  &=0
\end{align}
In \tfem\ a more general version of the ternary free energy is implemented, 
i.e.\
\begin{equation}
f_0 = \frac{1}{4} 
\left(
\alpha c_1^2 c_2^2 + \beta (c_1^2+c_2^2)\times(1-c_1-c_2)^2 - \gamma c_1c_2 
(1-c_1-c_2)
\right)
\end{equation}

A one-dimensional example of the ternary diffuse-interface model
can be run by
\begin{quote}
   \texttt{diffuse\_interface9 < input\_diffuse\_interface9.txt}
\end{quote}
and an two-dimensional example of the ternary diffuse-interface model
can be run by
\begin{quote}
   \texttt{diffuse\_interface10 < input\_diffuse\_interface10.txt}
\end{quote}

Another type of the ternary diffuse-interface model is also available in \tfem.
Detail of the model can be found in \cite{Boyer2006},
and only the relevant equations are summarized here.
The total free energy reads
\begin{equation}
f = f_0 + \frac{1}{2} \epsilon_1 | \nabla c_1 |^2
        + \frac{1}{2} \epsilon_2 | \nabla c_2 |^2
        + \frac{1}{2} \epsilon_3 | \nabla c_3 |^2
\end{equation}
where the homogeneous part $f_0$ is given by
\begin{equation}
f_0(c_1, c_2, c_3)
    = \frac{1}{2} \alpha c_1^2 \left( 1 - c_1 \right)^2
    + \frac{1}{2} \beta c_2^2 \left( 1 - c_2 \right)^2
    + \frac{1}{2} \gamma c_3^2 \left( 1 - c_3 \right)^2
    + 3 \Omega c_1^2 c_2^2 c_3^2
\end{equation}
The convection-diffusion equation for $c_i, (i = 1, 2, 3)$ reads
\begin{equation}
\pderiv{c_i}{t} + \vec{u} \cdot \nabla c_i =
\nabla \cdot \left( M_i \nabla \mu_i \right)
\end{equation}
The chemical potential $\mu_i, (i = 1, 2, 3)$ is given by
\begin{equation}
\mu_i = \pderiv{f_0}{c_i} - \epsilon_i \nabla^2 c_i + \lambda
\end{equation}
where $\lambda$ is a Lagrange multiplier
to ensure the constraint of $c_1 + c_2 + c_3 = 1$.
With a postulate of $M_1 \epsilon_1 = M_2 \epsilon_2 = M_3 \epsilon_3$,
the Lagrange multiplier can be defined by
\begin{equation}
\lambda = - \frac{1}{M_s} \left( M_1 \pderiv{f_0}{c_1}
        + M_2 \pderiv{f_0}{c_2}
        + M_3 \pderiv{f_0}{c_3} \right)
\end{equation}
where $M_s = M_1 + M_2 + M_3$.
Therefore, the chemical potential can be rewritten as
\begin{align}
\mu_1 = \frac{M_2}{M_s} \left( \pderiv{f_0}{c_1} - \pderiv{f_0}{c_2} \right)
      + \frac{M_3}{M_s} \left( \pderiv{f_0}{c_1} - \pderiv{f_0}{c_3} \right)
      - \epsilon_1 \nabla^2 c_1
\\
\mu_2 = \frac{M_3}{M_s} \left( \pderiv{f_0}{c_2} - \pderiv{f_0}{c_3} \right)
      + \frac{M_1}{M_s} \left( \pderiv{f_0}{c_2} - \pderiv{f_0}{c_1} \right)
      - \epsilon_2 \nabla^2 c_2
\\
\mu_3 = \frac{M_1}{M_s} \left( \pderiv{f_0}{c_3} - \pderiv{f_0}{c_1} \right)
      + \frac{M_2}{M_s} \left( \pderiv{f_0}{c_3} - \pderiv{f_0}{c_2} \right)
      - \epsilon_3 \nabla^2 c_3
\end{align}
It can be noted that $M_1 \mu_1 + M_2 \mu_2 + M_3 \mu_3 = 0$.
Finally, the system of equations for $\mu_1, \mu_2, c_1$ and $c_2$ reads
\begin{align}
& \pderiv{c_1}{t} + \vec{u} \cdot \nabla c_1 =
\nabla \cdot \left( M_1 \nabla \mu_1 \right)
\\
& \pderiv{c_2}{t} + \vec{u} \cdot \nabla c_2 =
\nabla \cdot \left( M_2 \nabla \mu_2 \right)
\\
& \mu_1 = \frac{M_2}{M_s} \left( \pderiv{f_0}{c_1} - \pderiv{f_0}{c_2} \right)
      + \frac{M_3}{M_s} \left( \pderiv{f_0}{c_1} - \pderiv{f_0}{c_3} \right)
      - \epsilon_1 \nabla^2 c_1
\\
& \mu_2 = \frac{M_3}{M_s} \left( \pderiv{f_0}{c_2} - \pderiv{f_0}{c_3} \right)
      + \frac{M_1}{M_s} \left( \pderiv{f_0}{c_2} - \pderiv{f_0}{c_1} \right)
      - \epsilon_2 \nabla^2 c_2
\end{align}
which is solved in combination with the Stokes equation
\begin{align}
0& = -\nabla p + \eta\nabla^2 {\vec u}
   + \rho (\mu_1\nabla c_1 + \mu_2\nabla c_2 + \blue{\mu_3 \nabla c_3})
\\
\nabla \cdot {\vec u}  &=0
\end{align}
Because $\nabla c_3 = - \nabla c_1 - \nabla c_2$ and
$\mu_3 = - \dfrac{M_1}{M_3} \mu_1 - \dfrac{M_2}{M_3} \mu_2$,
the momentum conservation equation actually implemented in \tfem\ is
\begin{equation}
0 = -\nabla p + \eta\nabla^2 {\vec u}
   + \rho \left( (\mu_1 + \blue{ \frac{M_1}{M_3} \mu_1
                               + \frac{M_2}{M_3} \mu_2 } ) \nabla c_1
               + (\mu_2 + \blue{ \frac{M_1}{M_3} \mu_1
                               + \frac{M_2}{M_3} \mu_2 } )\nabla c_2
          \right)
\end{equation}

The example file is \texttt{diffuse\_interface19.f90}, and
there are two different input files (\texttt{input\_diffuse\_interface19a.txt} and
\texttt{input\_diffuse\_interface19b.txt}).
A drop is initially located between two stratified fluids.
With the former input file, the three phases are fully symmetric and the drop will partially spread between the upper and lower fluid.
In the latter case, the upper fluid is totally spreading between the drop and lower fluid as shown in Fig. \ref{fig:ternary_spread}.

Note, the coefficients $M_2$ and $M_3$ should not be defined in the input file
but they are computed in the element routine via the relationship of
$M_1 \epsilon_1 = M_2 \epsilon_2 = M_3 \epsilon_3$.

\begin{figure}[htbp!]
   \begin{center}
   \includegraphics[width=0.3\textwidth]{ternary_spread_1}
   \includegraphics[width=0.3\textwidth]{ternary_spread_2}
   \includegraphics[width=0.3\textwidth]{ternary_spread_3}
   \end{center}
   \caption{Drop between two fluids:
            the upper fluid is totally spreading between
            the drop and lower fluid.
            $\alpha = 3$, $\beta = -1$, $\gamma = 3$, $\Omega = 7$,
            $M_1 = 2 \times 10^{-2}$, $\epsilon_1 = 12 \times 10^{-4}$,
            $\epsilon_2 = -4 \times 10^{-4}$,
            $\epsilon_3 = 12 \times 10^{-4}$,
            and time step $\Delta t = 5 \times 10^{-4}$.}
    \label{fig:ternary_spread}
\end{figure}

\subsection{Static contact angle}
\label{sec:contactangle}

(This text and examples in this section have been supplied by Nick Jaensson of the TU/e).

When two fluids are in contact with a solid, the angle the fluid-fluid 
interface makes with the solid interface is called the contact angle and is 
denoted by $\theta_\text{c}$, as is shown in Fig.~\ref{fig:droplet_on_solid}. 
The contact angle is related to the 
interfacial tensions/energies through Young's equation, given by
\begin{equation}
\gamma_{sb} = \gamma_{sa} + \gamma \cos{\theta_\text{c}},
\label{eq:youngs_eq}
\end{equation}
where $\gamma_{sb}$, $\gamma_{sa}$ and $\gamma$ are the interfacial 
tensions/energies 
between fluid $b$ and the solid, between fluid $a$ and the solid and between 
the two fluids, respectively.
\begin{figure}[htbp!]
   \begin{center}
   \includegraphics[scale=1.2]{droplet_on_solid}
   \end{center}
   \caption{Droplet of fluid $a$ , surrounded by fluid $b$, in contact with a 
      solid surface.}
    \label{fig:droplet_on_solid}
\end{figure}
In the diffuse-interface model, a contact angle that is not 
equal to $90^\circ$ can be imposed by adding a wall energy that is a function 
of $c$ \cite{Khatavkar2007}
\begin{equation}
F = \int_{\Omega} \rho f \, dV + \int_{\Gamma} \rho f_\text{w} \, dS,
\label{eq:total_free_energy}
\end{equation}
where $F$ is the total free energy and $f$ and $f_\text{w}$ are the 
(specific) free energies in the bulk and at the wall.
Furthermore, $\Gamma$ is the 
boundary where the contact angle boundary conditions is applied.
Using variational calculus, the wall energy in Eq.~\eqref{eq:total_free_energy} 
leads to the following boundary condition \cite{Beveridge1970} 
\begin{equation}
\kappa \frac{\partial c}{\partial n } + \frac{\partial 
f_\text{w}}{\partial c} 
= 0 \,\,\,\ \text{on} \,\, \Gamma, \label{eq:bc_contact_angle}
\end{equation}
where $n$ is the normal to the boundary $\Gamma$, pointing away from the 
fluids. Note, that $\rho$ is assumed constant and was divided out in 
Eq.~\eqref{eq:bc_contact_angle}.
For $f_\text{w}$, we use the following function, which was proposed by Jacqmin 
\cite{Jacqmin2000} 
\begin{equation}
f_\text{w} = \phi \left(c-\frac{c^3}{3} \right), \label{eq:fw}
\end{equation}
where $\phi$ is the wetting potential. We can relate $\phi$ to the contact 
angle by considering the change in wall energy $\Delta 
f_\text{w}$ across the 
fluid-fluid interface, as one moves along the solid boundary. Assume, for 
example, that the composition of fluid $a$ in Fig.~\ref{fig:droplet_on_solid} 
is given by $-c_\text{B}$, while the composition of fluid $b$ is given by 
$c_\text{B}$. Moving along the solid boundary from fluid $a$ to fluid $b$, the 
change in wall energy is given by
\begin{equation}
\Delta f_\text{w}=f_\text{w}(c_\text{B}) - f_\text{w}(-c_\text{B}).
\label{eq:Delta_f1}
\end{equation}
From 
Eq.~\eqref{eq:youngs_eq} follows that the difference between $f_{sa}$ and 
$f_{sb}$ (times $\rho$ since $f_\text{w}$ is the specific wall free energy) 
is given by $\gamma \cos \theta_\text{c}$, therefore 
\begin{equation}
\rho \Delta f_\text{w}=\gamma \cos \theta_\text{c}. \label{eq:Delta_f2}
\end{equation}
Equating Eqs.~\eqref{eq:Delta_f1} and \eqref{eq:Delta_f2} and using the definition 
of $f_w$ given by Eq.~\eqref{eq:fw} yields an expression for the wetting 
potential in terms of the contact angle and (known) diffuse-interface 
coefficients
\begin{equation}
\phi = \frac{\cos \theta_\text{c}}{2(c_\text{B} -c_\text{B}^3/3 )} 
\frac{\gamma}{\rho}.
\label{eq:wetting_potential}
\end{equation}
The approximation for the surface tension of a planar interface, denoted by
$\gamma$ and given in Eq.~\eqref{eq:di_surface_tension} can be used to
 determine $\phi$ in Eq.~\eqref{eq:wetting_potential}.

\subsubsection{Numerical implementation}
Rewriting the weak from of Eq.~\eqref{eq:di4}, but now without the boundary 
condition $\partial c / \partial n = 0$ yields
\begin{equation}
  (s, \alpha c-\beta c^3) -   \kappa(\nabla s,\nabla c) + ( s, \kappa
  \frac{\partial c}{\partial n})_{\Gamma}
  +(s,\mu) =0 \qquad\text{for all } s.
\end{equation}
The boundary condition given by Eq.~\eqref{eq:bc_contact_angle} can now
be directly substituted into the third term on the left hand side, which is 
then given by
\begin{align}
  ( s, \kappa  \frac{\partial c}{\partial  n} )_{\Gamma} 
  &= -( s, \frac{\partial f_\text{w}}{\partial c} )_{\Gamma} \\
  &= -\left( s, \phi (1-c^2) \right)_{\Gamma} \label{eq:weak_di_bc}
\end{align}
Note, that Eq.~\eqref{eq:weak_di_bc} is non-linear and can be solved by 
a Picard iteration
\begin{align}
 ...+\left(s,\phi c_i\right)_{\Gamma}c_{i+1} = 
 ...+\left(s,\phi\right)_{\Gamma},
\end{align}
where the dots mean that this equation is added to Eq.~\eqref{eq:di_picard}. 
When using a Newton-Raphson iteration, Eq.~\eqref{eq:weak_di_bc} becomes
\begin{align}
 ...+\left(s,2\phi c_i\right)_{\Gamma}c_{i+1} = 
 ...+\left(s,\phi(1+c_i^2)\right)_{\Gamma},
\end{align}
where the dots mean that this equation is added to Eq.~\eqref{eq:di_newton}.

In the example file \texttt{diffuse\_interface13.f90} a droplet of fluid 
with $c=-1$, surrounded by fluid with $c=1$, lying on the bottom of a closed 
cylindrical 
container is simulated. The problem is solved using the assumption of 
axisymmetry.
The contact angle boundary condition is imposed on the bottom of the cylinder, 
where the contact angle measured through the droplet is given in 
degrees in \texttt{theta\_c}. In this example, first-order time integration and 
a mesh with a uniform element size is used. Furthermore, the diffuse-interface  
equations are solved using a Newton iteration.

More example files are \texttt{diffuse\_interface14.f90}, where again a droplet 
on the bottom of a cylindrical container is solved using the axisymmetry 
assumption, and \texttt{diffuse\_interface15.f90} where a droplet on the bottom 
of a closed cylindrical container is simulated in 3D. Note, that in 
\texttt{diffuse\_interface14.f90} and \texttt{diffuse\_interface15.f90} the 
time integration is second-order. Furthermore, adaptive meshing is used, as 
will be explained in the following. In 3D, solving the system of equations 
can take a long time using a sequential solver. In that case, 
\texttt{diffuse\_interface16.f90} can be used, which is similar to 
\texttt{diffuse\_interface15.f90}, but now using Pardiso as the (parallel) 
solver.

\subsubsection{Adaptive meshing}
Gmsh (see Sec.~\ref{sec:readgmshmesh}) is used to generate meshes that are
refined in the interface 
region, but coarse far way. Gmsh can use ``refinement fields'' to
locally control the element size. Refinement fields can be
defined using refinement points: 
points that are not part of the mesh (geometrical points), that 
define an element size $h_\text{fine}$ in the circle/sphere with 
radius $d_\text{min}$, as is shown in Fig.~\ref{fig:refinement_point}. Outside 
the circle/sphere with radius $d_\text{max}$, element size $h_\text{coarse}$ 
is used, with a transitional region in $d_\text{max}-d_\text{min}$. Each 
refinement field can contains multiple refinement points (with the same $h_\text{fine}$, $h_\text{coarse}$, $d_\text{min}$ and $d_\text{max}$). Furthermore,
multiple refinement fields can be used simultaneously (at every point in space, the smallest target element size is used when meshing).   To make sure
Gmsh does not specify the elementsize by itself, it
advised to use the option 
\texttt{Mesh.CharacteristicLengthExtendFromBoundary = 0}. For the 
diffuse-interface simulations, the refinement points have to initially be given 
manually, but during the simulation, refinement points can be generated using 
the $c$-field. The approach taken here is to generate a list of
refinement points by adding a refinement point in each element 
that 
crosses the isoline $c=0$. The coordinates of the refinement point are given by 
the coordinates of the nodal point with the smallest $|c|$, i.e. the nodal 
point that is closest to the interface. 
\begin{figure}[htbp!]
   \begin{center}
   \includegraphics[scale=1.2]{refinement_point}
   \end{center}
   \caption{A refinement point (with coordinates $(x,y)$) can be used in Gmsh 
   to enforce a 
   element size $h_\text{fine}$ in the circle with radius $d_\text{min}$ 
   around the refinement point, while a larger elements size $h_\text{coarse}$ 
   is used outside the circle with radius $d_\text{max}$.}
    \label{fig:refinement_point}
\end{figure}
Since the interface might move out of the refined region, a criterion is needed 
to invoke remeshing when the element size around the interface becomes too 
large. An approach similar
to \cite{Yue2006} is used, where a 
ribbon of elements around the interface is defined which are allowed to contain 
the interface (given by the isoline $c=0$). As soon as this isoline moves into 
a non-allowed element, remeshing is invoked. The distance to the isoline $c=0$ 
could 
be computed to determine which element is an allowed element, but it is more 
convenient to use the nodal values of $c$, since these are readily available. 
Allowed elements are given by elements for which \emph{any} nodal value of 
$|c| < c^*$, where  $c^*$ is a threshold value. A graphical representation of 
this criterion is shown in Fig.~\ref{fig:allowed_elements}. Examples for using 
the refinement fields for adaptive meshing are given in
\texttt{diffuse\_interface14.f90} and \texttt{diffuse\_interface15.f90}. 
The code to include the refinement points when meshing with Gmsh can be found
in the corresponding \texttt{.igo} files.
\begin{figure}[htbp!]
   \begin{center}
   \includegraphics[scale=1.2]{allowed_elements}
   \end{center}
   \caption{Elements for which any of the nodal values of $|c|$ are smaller 
   than a treshold value $c^*$ are added to the list of allowed 
   elements (shown in green). When the interface (defined by the line $c=0$) 
   enters a non-allowed element (shown in red) remeshing is invoked.}
    \label{fig:allowed_elements}
\end{figure}

\subsection{Fully implicit time discretization}
The time discretization schemes shown in Section \ref{sec:di_expl} take the 
surface tension explicitly into account in the momentum balance. At high 
surface tension (or low capillary number), this can lead to instabilities of 
the interface when the time step is taken too large. As is shown in 
\cite{Aland2014}, a fully implicit scheme leads to large increase in stability 
and allows 
much larger time steps, while maintaining a stable interface. The fully 
implicit scheme is obtained by solving for $\vek u$, $g$, $c$ and $\mu$ in one 
system. The terms $\rho \kappa (|\nabla c|^2 \ten I 
- \nabla c \nabla c)$ (when using the stress form), $\rho \mu \nabla c$ (when 
using potential form 1), $\rho c \nabla \mu$ (when 
using potential form 2) and $\vek u \cdot 
\nabla c$ are quadratic for which we can use a Picard iteration (note that 
$\mu$ is chosen as the implicit unknown in the potential forms, which has the 
largest stabilizing effect \cite{Aland2014}).
Using an implicit Euler scheme for the time derivative of $c$, we obtain the 
weak form for the stress form:
find $\vek u_{i+1}$, $[g_0]_{i+1}$, $c_{i+1}$ and $\mu_{i+1}$ such that
\begin{align*}
\big(\ten D_v,2\eta \ten D_{i+1} \big) - (\nabla\cdot\vek v,[g_0]_{i+1})
+\blue{(\ten D_v, \rho \kappa(\nabla c_i \cdot \nabla c_{i+1} 
	\ten I } \blue{- \nabla c_i \nabla c_{i+1}))}&= \\
(\vek v,\vek t_{\text{N}})_{\Gamma_{\text{N}}}+&\;(\vek v, \vek f)
&\qquad&\text{for all }\vek v\\
(q,\nabla\cdot\vek u_{i+1})&=0&\qquad&\text{for all }q\\
\big(r,\frac{c_{i+1}-c^n}{\Delta t}+\blue{\vek u_i}\cdot\nabla c_{i+1}\big)
+M(\nabla r,\nabla \mu_{i+1})&=0&\qquad&\text{for all } r\\
\big(s, (\alpha -\beta c_{i}^2)c_{i+1}\big)-\kappa(\nabla s,\nabla c_{i+1}) 
+(s,\mu_{i+1})&=0&\qquad&\text{for all }s
\end{align*}
where the  
terms that arise due to implicit time discretization are marked in blue. Using 
potential form 1, we obtain the weak form:
find $\vek u_{i+1}$, $[g_1]_{i+1}$, $c_{i+1}$ and $\mu_{i+1}$ such that
\begin{alignat*}{6}
\big(\ten D_v,2\eta \ten D_{i+1} \big) -(\nabla\cdot\vek v,[g_1]_{i+1})
-\blue{( \vek v, \rho \mu_{i+1}\nabla c_{i})}&=(\vek v, \vek f)
&\qquad&\text{for all }\vek v\\
(q,\nabla\cdot\vek u_{i+1})&=0&\qquad&\text{for all }q \\
\big(r,\frac{c_{i+1}-c^n}{\Delta t}+\blue{\vek u_i}\cdot\nabla c_{i+1}\big)
+M(\nabla r,\nabla \mu_{i+1})&=0 &\qquad&\text{for all } r\\
\big(s, (\alpha -\beta c_{i}^2)c_{i+1}\big)-\kappa(\nabla s,\nabla c_{i+1}) 
+(s,\mu_{i+1})&=0&\qquad&\text{for all } s
\end{alignat*}
Using potential form 2, we obtain the weak form:
find $\vek u_{i+1}$, $[g_2]_{i+1}$, $c_{i+1}$ and $\mu_{i+1}$ such that
\begin{alignat*}{6}
\big(\ten D_v,2\eta \ten D_{i+1} \big) -(\nabla\cdot\vek v,[g_2]_{i+1})
+\blue{( \vek v, \rho c_{i}\nabla \mu_{i+1})}&=(\vek v, \vek f)
&\qquad&\text{for all }\vek v\\
(q,\nabla\cdot\vek u_{i+1})&=0&\qquad&\text{for all }q \\
\big(r,\frac{c_{i+1}-c^n}{\Delta t}+\blue{\vek u_i}\cdot\nabla c_{i+1}\big)
+M(\nabla r,\nabla \mu_{i+1})&=0 &\qquad&\text{for all } r\\
\big(s, (\alpha -\beta c_{i}^2)c_{i+1}\big)-\kappa(\nabla s,\nabla c_{i+1}) 
+(s,\mu_{i+1})&=0&\qquad&\text{for all } s
\end{alignat*}

Alternatively, we can handle the quadratic terms with a Newton-Raphson
iteration (similar to Section \ref{sec:nonlin}):
\begin{alignat*}{3}
\rho \kappa (|\nabla c|^2 \ten I -\nabla c \nabla c) &\approx 
&& 2\rho \kappa \nabla c_i \cdot \nabla c_{i+1} \ten I -\rho \kappa 
|\nabla c_i|^2 \ten I-\rho \kappa \nabla c_{i+1} \nabla c_i -\rho \kappa 
\nabla c_i \nabla c_{i+1}\\
&\quad\; &&+ \rho \kappa \nabla c_i \nabla c_i\\
\rho \mu \nabla c &\approx &&\rho \mu_{i+1} \nabla c_{i}+\rho \mu_{i} 
\nabla c_{i+1} - \rho \mu_{i} \nabla c_{i} \\
-\rho c \nabla \mu &\approx &&-\rho c_{i+1} \nabla \mu_{i}-\rho c_{i} 
\nabla \mu_{i+1} + \rho c_{i} \nabla \mu_{i} \\
\vek u \cdot \nabla c &\approx &&\vek u_{i+1} \cdot \nabla c_{i} + \vek u_{i} 
\cdot \nabla c_{i+1} - \vek u_{i} \cdot \nabla c_{i} 
\end{alignat*}
Substitution of these terms (and using a Newton-Raphson iteration for the 
$c^3$-term, similar to Eq.~\eqref{eq:di_newton}) 
leads to the weak form for the stress form:
find $\vek u_{i+1}$, $[g_0]_{i+1}$, $c_{i+1}$ and $\mu_{i+1}$ such that
\begin{align*}
\big(\ten D_v,2\eta \ten D_{i+1} \big) - (\nabla\cdot\vek v,[g_0]_{i+1})
+\blue{(\ten D_v, \rho \kappa(2\nabla c_i \cdot \nabla c_{i+1} \ten I }
\blue{-\nabla c_i \nabla c_{i+1}}&\blue{-\nabla c_{i+1} \nabla c_{i}))}=\\
(\vek v,\vek t_{\text{N}})_{\Gamma_{\text{N}}}+(\vek v, \vek f) 
+\blue{(\ten D_v, \rho \kappa (|c_i|^2\ten I}&\blue{-\nabla c_i 
	\nabla c_i))}\\
(q,\nabla\cdot\vek u_{i+1})&=0\\
\big(r,\frac{c_{i+1}-c^n}{\Delta t}+\blue{\vek u_i}\cdot\nabla c_{i+1}\big)
+\blue{\big(r, \vek u_{i+1} \cdot \nabla c_i \big)}+M(\nabla 
r,\nabla \mu_{i+1})&=\blue{(r, \vek u_i \cdot \nabla c_i)} \\
\big(s, (\alpha -3\beta c_{i}^2)c_{i+1}\big)-\kappa(\nabla s,\nabla c_{i+1})
+(s,\mu_{i+1})&=-\big(s,2\beta c_{i}^3 \big)
\end{align*}
for all $\vek v$, $q$, $r$ and $s$, respectively. Using potential form 1, we 
obtain the weak form:
find $\vek u_{i+1}$, $[g_1]_{i+1}$, $c_{i+1}$ and $\mu_{i+1}$ such that
\begin{alignat*}{6}
\big(\ten D_v,2\eta \ten D_{i+1} \big)-(\nabla\cdot\vek v,[g_1]_{i+1})
-\blue{(\vek v,\rho \mu_{i+1}\nabla c_{i})}-\blue{( \vek v,\rho \mu_{i}
	\nabla c_{i+1})}&=(\vek v,\vek f)-\blue{(\vek v,\rho \mu_{i}\nabla c_{i})}\\
(q,\nabla\cdot\vek u_{i+1})&=0\\
\big(r,\frac{c_{i+1}-c^n}{\Delta t}+\blue{\vek u_i}\cdot\nabla c_{i+1}\big)
+\blue{\big(r, \vek u_{i+1} \cdot \nabla c_i \big)}+M(\nabla r,\nabla 
\mu_{i+1})&=\blue{(r, \vek u_i \cdot \nabla c_i)} \\
\big(s, (\alpha -3\beta c_{i}^2)c_{i+1}\big)-\kappa(\nabla s,\nabla c_{i+1}) 
+(s,\mu_{i+1})&=-\big(s, 2\beta c_{i}^3 \big)
\end{alignat*}
for all $\vek v$, $q$, $r$ and $s$, respectively. Using potential form 2, we 
obtain the weak form:
find $\vek u_{i+1}$, $[g_2]_{i+1}$, $c_{i+1}$ and $\mu_{i+1}$ such that
\begin{alignat*}{6}
\big(\ten D_v,2\eta \ten D_{i+1} \big) -(\nabla\cdot\vek v,[g_2]_{i+1})   
+\blue{( \vek v, \rho c_{i}\nabla \mu_{i+1})}+\blue{( \vek v, \rho c_{i+1} 
	\nabla \mu_{i})}&=(\vek v,\vek f)+\blue{(\vek v,\rho c_{i}\nabla \mu_{i})}\\
(q,\nabla\cdot\vek u_{i+1})&=0\\
\big(r,\frac{c_{i+1}-c^n}{\Delta t}+\blue{\vek u_i}\cdot\nabla c_{i+1}\big) 
+\blue{\big(r, \vek u_{i+1} \cdot \nabla c_i \big)}+M(\nabla r,\nabla
\mu_{i+1})&=\blue{(r, \vek u_i \cdot \nabla c_i)} \\
\big(s, (\alpha -3\beta c_{i}^2)c_{i+1}\big)-\kappa(\nabla s,\nabla c_{i+1}) 
+(s,\mu_{i+1})&=-\big(s, 2\beta c_{i}^3 \big)
\end{alignat*}
for all $\vek v$, $q$, $r$ and $s$, respectively.

An example of the diffuse-interface model with fully implicit time 
discretization is \texttt{diffuse\_interface18} (the same problem but with 
explicit/implicit time discretization is  \texttt{diffuse\_interface17}). In 
these examples, an axisymmetric droplet (initially spherical) is simulated in a 
hyperbolic flow. Depending on the capillary number 
$\text{Ca}=\eta_\text{d}\dot{\epsilon} r / \gamma$ (where 
$\eta_\text{d}$ is the droplet viscosity, $\dot{\epsilon}$ is the elongation 
rate, $r$ is the initial drop radius and $\gamma$ is the surface tension), the 
droplet will deform and possibly break up.

Remarks:
\begin{enumerate}
	\item Potential form 1, potential form 2 and the stress 
	form are implemented in both examples, and can be used to add the
        effect of surface tension to the momentum balance.
	\item It is possible to have a different viscosities for 
	the droplet and the ambient fluid. A linear mixing rule is used:
	\[
	\eta(c) = \eta_1 \frac{c+c_\text{B}}{2c_\text{B}} - \eta_2   
	\frac{c-c_\text{B}}{2c_\text{B}}
	\]
	where $c_B$ is the concentration in the bulk and $\eta_1$ and $\eta_2$ 
        are the viscosities for $+c_\text{B}$ and $-c_\text{B}$, respectively. 
	For the implicit scheme, a second order prediction 
        $\hat{c}=2c_n-c_{n-1}$ is used to determine 
	$\eta(c)$, which can be combined with a Newton-Raphson iteration.
	Alternatively, the variable viscosity can be included in the iteration,
        but this limits the  
	iteration to Picard (the approximation $\eta(c)\ten D \approx 
	\eta(c_i)\ten D_{i+1}$ is used). By default, the second-order prediction
        is used in \texttt{diffuse\_interface18}, but the implicit viscosity 
        iteration is implemented as well.
	\item Since in these examples a variable viscosity is used, the system 
	matrix in the momentum 
	balance changes between time steps and iteration steps and has to be 
	rebuild (also the LU decomposition can not be reused).
\end{enumerate}
